\documentclass[letterpaper,12pt]{article}
\usepackage[paperheight=27.94cm,paperwidth=21.59cm,bindingoffset=0in,left=3cm,right=2.0cm,top=3.5cm,bottom=2.5cm,headheight=15pt,headsep=1.0\baselineskip]{geometry}
\usepackage{graphicx}
\usepackage{amsmath,amssymb}
\usepackage[spanish,es-tabla]{babel}
\usepackage{booktabs}
\usepackage{tabularx}
\usepackage{float}
\usepackage{cite}
\usepackage[utf8]{inputenc}
\usepackage{fancyhdr}
\usepackage{xurl}
\usepackage[hidelinks]{hyperref}
\usepackage{listings}
\usepackage{xcolor}

% Configuración del formato para el código VHDL
\lstset{
    language=VHDL,
    basicstyle=\ttfamily\footnotesize,
    keywordstyle=\color{blue}\bfseries,
    commentstyle=\color{green!60!black},
    stringstyle=\color{red},
    numbers=left,
    numberstyle=\tiny\color{gray},
    stepnumber=1,
    numbersep=10pt,
    backgroundcolor=\color{black!5},
    frame=single,
    rulecolor=\color{black!30},
    breaklines=true,
    breakatwhitespace=false,
    showstringspaces=false,
    captionpos=b,
    inputencoding=utf8,
    extendedchars=true,
    literate=
        {á}{{\'a}}1 {é}{{\'e}}1 {í}{{\'i}}1 {ó}{{\'o}}1 {ú}{{\'u}}1
        {ñ}{{\~n}}1 {Á}{{\'A}}1 {É}{{\'E}}1 {Í}{{\'I}}1 {Ó}{{\'O}}1
        {Ú}{{\'U}}1 {Ñ}{{\~N}}1 {¡}{{\textexclamdown}}1 {¿}{{\textquestiondown}}1
}

\renewcommand{\tablename}{Tabla}
\pagestyle{fancy}
\fancyhf{}
\rhead{\thepage}
\lhead{Laboratorio 3: Máquina de estados finitos para control de acceso digital}

\begin{document}

% ==============================
% 1. CARÁTULA
% ==============================
\begin{titlepage}
    \centering
    \includegraphics[width=0.55\textwidth]{Logo_Universidad_Diego_Portales.png}\\[1.2 cm]
    
    \textsc{\large Facultad de Ingeniería y Ciencias}\\[0.2 cm]
    \textsc{\large Escuela de Informática y Telecomunicaciones}\\[1.5 cm]
    
    \textsc{\Large Arquitectura y Organización de Computadores}\\[0.5 cm]
    
    \rule{\linewidth}{0.5 mm} \\[0.4 cm]
    {\huge \bfseries Laboratorio 4:\\[0.3cm] \Large Temporizador regresivo y contador de eventos en FPGA}\\
    \rule{\linewidth}{0.5 mm} \\[2.0 cm]
    
    \large \emph{Integrantes:}\\[0.3 cm]
    {\Large Benjamín Oyarzún}\\[0.2cm]
    {\Large Lucas Herrada}\\[0.2cm]
    {\Large Joakin Mac-Auliffe}

    \vfill
    {\large Diciembre de 2026}
\end{titlepage}

\newpage
\tableofcontents
\newpage

% ==============================
% 2. INTRODUCCIÓN
% ==============================
\section{Introducción}
.

% ==============================
% 3. OBJETIVOS
% ==============================
\section{Objetivos}


% ==============================
% 4. FUNDAMENTOS TEÓRICOS
% ==============================
\section{Fundamentos teóricos}


% ==============================
% 5. DESCRIPCIÓN FUNCIONAL
% ==============================
\section{Descripción funcional}

% ==============================
% 6. DIAGRAMA DE BLOQUES
% ==============================
\section{Diagrama de bloques}


% ==============================
% 7. DIAGRAMA DE ESTADOS
% ==============================
\section{Diagrama de estados}

% ==============================
% 8. TABLA DE TRANSICIÓN
% ==============================
\section{Tabla de transición}

\section{Código VHDL}

% ==============================
% 10. TESTBENCH
% ==============================
\section{Testbench}


% ==============================
% 11. SIMULACIÓN EN VIVADO
% ==============================
\section{Simulación en Vivado}



% ==============================
% 12. ARCHIVO DE RESTRICCIONES
% ==============================
\section{Archivo de restricciones}


\section{Resultados obtenidos}

% ==============================
% 13. ANÁLISIS DE ERRORES
% ==============================
\section{Análisis de errores}

% ==============================
% 14. RESPUESTAS A PREGUNTAS
% ==============================
\section{Respuestas a preguntas}

% ==============================
% 15. CONCLUSIONES
% ==============================
\section{Conclusiones}

\newpage
% ==============================
% REFERENCIAS
% ==============================
\addcontentsline{toc}{section}{Referencias}
\begin{thebibliography}{9}

\bibitem{nexys4ddr}
Digilent Inc., \emph{Nexys 4 DDR FPGA Board Reference Manual}, DOC\#: 502-292. [En línea]. Disponible en: \url{https://digilent.com/reference/programmable-logic/nexys-4-ddr/reference-manual}

\bibitem{artix7}
Xilinx, Inc., \emph{7 Series FPGAs Data Sheet: Overview}, DS180.

\bibitem{ieee1076}
IEEE, \emph{IEEE Standard VHDL Language Reference Manual}, IEEE Std 1076-2008, 2009.

\bibitem{ieee1164}
IEEE, \emph{IEEE Standard Multivalue Logic System for VHDL Model Interoperability (Std\_logic\_1164)}, IEEE Std 1164-1993, 1993.

\bibitem{ieee1076_3}
IEEE, \emph{IEEE Standard VHDL Synthesis Packages}, IEEE Std 1076.3-1997, 1997.

\bibitem{harris}
D. M. Harris y S. L. Harris, \emph{Digital Design and Computer Architecture}, 2.ª ed. Morgan Kaufmann, 2012.

\bibitem{pedroni}
V. A. Pedroni, \emph{Circuit Design and Simulation with VHDL}, 2.ª ed. MIT Press, 2010.

\end{thebibliography}

\newpage
% ==============================
% 7. ANEXOS
% ==============================
\section{Anexos}

\subsection{Repositorio del proyecto}

\begin{itemize}
    \item \url{https://github.com/JoakinMacAuliffe/Lab4_Arqui}
\end{itemize}

\subsection{Código fuente VHDL}

\subsubsection{\texttt{detector\_flanco.vhd} (Módulo sincronizador y detector de flancos)}

\subsubsection{\texttt{lab03\_top.vhd} (El módulo principal que integra el diseño)}
\begin{lstlisting}[language=VHDL, caption={lab03\_top.vhd}]

\end{lstlisting}

\subsection{Testbench (Banco de pruebas)}

\subsection{Archivo XDC (Restricciones)}

\subsubsection{\texttt{lab03.xdc} (La asignación física de pines para la FPGA)}
\begin{lstlisting}[language=tcl, caption={lab03.xdc}]

\end{lstlisting}

\end{document}